\chapter{Introduction}\label{chap:introduction}

Partial Order Causal Link (POCL) planning is a solving technique in planning which maintains the principle of least commitment, as explained by~\cite{Weld1994LeastCommitmentPlanning}. While state-based planning approaches search the space of states, requiring commitment to a specific ordered sequence of actions which reaches that state, POCL planning searches the space of partial plans, where actions only have a partial set of ordering constraints, and so no such commitments need to be made. As explained by~\cite{Minton1994TotalVsPartialPlanning}, a single partial plan can hence represent up to a factorial number of action sequences, and hence POCL search trees can be exponentially smaller than their total-order counterparts, which~\cite{Weld2011NonlinearPlanningCommentary} notes was used as a justification for the intuitive superiority of partial-order planning until the mid 90s. However,~\cite{Minton1994TotalVsPartialPlanning} also found that this assumption of superiority required question, since `the superiority of partial-order planning \textit{can} depend critically upon the search strategy and search heuristics employed'. Total-order state based planning allows for intuitive, easy to compute, and powerful heuristics and search strategies, whereas partial-order planning struggled to generate any comparable heuristics, and this began to shift planning research toward total-order planning. Part of this shift was the development of Graphplan by~\cite{Blum1995Graphplan}, which took a fundamentally different approach to traditional partial-order planners, and outperformed UCPOP, the leading partial-order planner of the time created by~\cite{Penberthy1992UCPOP}. This shift continued with UNPOP, a state-based solver created by~\cite{McDermott1996UNPOP} which automatically generated domain-independent heuristics that allowed total-order state-based planning, which Graphplan had not fully embraced as it still produced partial-order outputs, to become competitive with the best planners of the time.\\


Today, state-based planning is still the dominant technique, with the FF planning system designed by~\cite{Hoffmann2001FFPlanning}, and then the Fast Downward planning system designed by~\cite{Helmert2006FastDownward}, significantly outperforming partial-order planners of the time such as VHPOP designed by~\cite{Younes2003VHPOPAndAdditiveHeuristic}. However, attempts at proving that the complexity of partial-order planning is higher than total-order planning has yielded fewer results. Both partial-order planning and total-order planning are PSPACE-complete problems, even in the simplest form as proven by~\cite{Bylander1994ComplexityOfSTRIPS}, with~\cite{Bercher2021POCLComplexities} confirming that this is also the case for POCL planning, and so any attempt to differentiate their complexity must investigate differences with the PSPACE-complete complexity class. Using sub-classes of PSPACE-completeness, such as those proposed by~\cite{Backstrom2011PSPACECompletePlanning}, still does not find partial-order planning to be higher in complexity than total-order planning. Instead, it is possible that heuristic efficiency could be at fault for the current poor performance of partial-order planners, a possibility which~\cite{Balcazar1996ImplicitGraphs} notes is important for all heuristic search algorithms on graphs.~\cite{Bercher2021POCLComplexities} finds that a common heuristic, the delete-relaxation heuristic, is NP-complete in POCL planning, compared to polytime for state-based planning, which provides further evidence for heuristic effectiveness and efficiency being a key drawback of POCL planning.\\


The input for heuristics in partial-order planning is more complex than state-based planning. In state-based planning, the current search progress can be fed to a heuristic algorithm as a constant size list of variables and their values. In POCL planning, this search progress contains a list of steps, which can contain repeated actions, a set of orderings between those steps, which can grow quadratically with the number of steps, and a set of causal links between effects and preconditions of these steps, which grows linearly with the number of steps. In addition, the size of each search node grows with search depth in POCL planning as more steps, orderings, and causal links are added, whereas it remains constant in state-based planning as the total number of existing propositional variables never changes. This problem has been recognised when creating heuristics for partial-order planning, and the two major well-informed domain-independent heuristics for POCL planning, the \textit{Relax heuristic} created by~\cite{Nguyen2001UCPOPAndRelaxHeuristic}, and the \textit{Additive heuristic} created by~\cite{Younes2003VHPOPAndAdditiveHeuristic}, solve it by including limited amounts of partial-plan information within their input. The additive heuristic ignores causal links and delete effects of actions within the current partial-plan, and the relax heuristic additionally ignores ordering constraints between actions, but doesn't make the assumption of subgoal independence, accounting for positive interactions between subgoals similarly to the FF heuristic by~\cite{Hoffmann2001FFPlanning}. This reduction of search node components to consider when generating heuristic estimates allows both of these heuristics to run in polytime.\\


The recognition of heuristic input complexity having a significant effect on planning efficiency causes a major issue in POCL planning. While ignoring most partial-plan information when calculating heuristics has led to the creation of informed heuristics that are very fast to compute, such as the Relax and Additive heuristics, these heuristics are significantly less informed than if they accounted for all of the information within the partial-plan they are calculated on. Causal links in particular contain a large amount of information, as they commit to causal relationships between actions, specifying that an effect of one action is required to protect the precondition of a later action, preventing actions that delete the linked proposition from being inserted between the two linked actions. This induces a constraint on where actions can be included within a plan, and hence implicitly prunes a large number of possible action sequences. The delete effects of actions within a partial plan also have a pruning effect, as they delete propositions which could have otherwise been used as preconditions for actions in places where those actions realistically couldn't appear if delete effects were included. Hence, causal links and delete effects encode a large amount of the search progress to reach a certain node, and by ignoring them within a heuristic, that heuristic is essentially being calculated for a significantly more abstract partial plan. This abstraction might be similar to going back many layers of depth in the search tree, and could be equivalent to many other explored search nodes, meaning that that these heuristics cannot differentiate between important differences in search nodes. In any other planning paradigm, running such a heuristic would likely be considered absurd, due to the fact that it is sacrificing large amounts of already computed search progress in order to save computation time in further expansion of the search tree.\\


We present new informed and admissible heuristics which take into account these often ignored components of a partial-plan search node. It has been shown by~\cite{Bercher2021POCLComplexities} that a delete-relaxed heuristic which respects ordering constraints; as well as either current causal links, delete effects, or both; will have an underlying decision problem which is NP-complete. However, as established previously, each POCL node can represent up to a factorial amount of action sequences, and so it is very possible that having a much more informed heuristic will unlock much of these benefits to partial-order planning that previous heuristics couldn't. Hence, an NP-complete heuristic might turn out to be worthwhile given the benefits to the search process. These heuristics are particularly promising in their application to satisfiability: determining whether a planning problem instance is solvable. State-based planners can struggle when presented with infeasible planning problems, as their smaller amount of information per search node compared to POCL planning can lead to a search tree that grows faster than infeasible branches can be pruned. This is also the case with the Relax and Additive heuristics for POCL planning, since they don't consider the large amount of additional information available within POCL search nodes. We propose that an informed POCL heuristic that takes into account all details of the much more complex structure of a POCL node could determine infeasibility for factorial numbers of action sequences in a single NP-complete calculation, possibly allowing it to prune branches faster than the search tree can grow in cases where state-based planning cannot, providing a potential benefit for determining infeasibility in planning.